% !TeX root = ../../tesi.tex

\section{Clientela e modalità di ingaggio}
\label{sec:clientela}

Halue lavora soltanto con aziende private, per lo più di medie e grandi dimensioni, e non con la
pubblica amministrazione. Sono marchi strutturati della moda, dello sport, del commercio al
dettaglio e della cosmetica professionale, con fatturati che vanno da decine a diverse centinaia di
milioni di euro all’anno, e vendono \emph{online} su più mercati, con lingue e valute diverse. La
maggior parte dei clienti è italiana, ma ci sono anche marchi esteri.

Vendere su più mercati moltiplica il lavoro tecnico. Ogni mercato ha la sua lingua, la sua valuta,
i suoi prezzi e le sue tasse, spesso anche un magazzino proprio, e ogni negozio deve restare
allineato con il sistema gestionale dell’azienda. Per questo molti clienti di Halue usano
\gloss{shopify-plus}{Shopify Plus}, la proposta di Shopify per le aziende di grandi dimensioni, che
aggiunge al negozio, tra l’altro, funzioni dedicate alla vendita
B2B\textsuperscript{\ref{rif:shopify-plus}}. Alcuni clienti hanno decine di negozi attivi, uno per
mercato o per canale, ciascuno con il proprio catalogo e il proprio magazzino.

Halue è \emph{Shopify Plus Partner} dal 2023. Il programma partner di Shopify assegna alle agenzie
un livello in base ai negozi che hanno avviato o trasferito sulla piattaforma e al volume di vendite
di quei negozi. Il livello dà all’agenzia visibilità nell’elenco ufficiale dei partner, che i
commercianti consultano per scegliere a chi affidare il proprio negozio. Nel suo profilo Halue
indica come specializzazioni, tra le altre, lo sviluppo di applicazioni e API su misura, la
personalizzazione del processo di pagamento e degli account dei clienti e l’ottimizzazione per i
motori di ricerca\textsuperscript{\ref{rif:halue-partner}}.

Le richieste dei clienti rientrano in quattro tipi:

\begin{itemize}
	\item la migrazione di un negozio verso Shopify da altre piattaforme, come Salesforce Commerce
	Cloud, WooCommerce e Magento. Oltre a catalogo, clienti e ordini, la migrazione trasferisce le
	integrazioni con i sistemi esterni e conserva gli indirizzi delle pagine, perché i collegamenti
	esistenti e il posizionamento nei motori di ricerca non vadano persi;
	\item il \gloss{redesign}{\emph{redesign}} di un negozio esistente;
	\item lo sviluppo di interfacce su misura, con \gloss{landing-page}{\emph{landing page}}
	dinamiche e configuratori 3D;
	\item l’espansione su nuovi mercati e la manutenzione periodica dei negozi in produzione, con
	aggiornamenti di configurazione, correzioni ed evoluzioni.
\end{itemize}

Durante il mio \emph{stage} le manutenzioni e gli aggiornamenti erano il tema più frequente nelle
riunioni quotidiane del \emph{team}.

I progetti arrivano soprattutto per passaparola, da clienti e contatti che segnalano l’azienda, e
in parte da gare e bandi. Le formule contrattuali sono tre, ciascuna adatta a un tipo di lavoro:

\begin{itemize}
	\item a corpo, con un prezzo fissato in anticipo, per i progetti dal perimetro ben definito.
	Questa formula richiede una stima accurata prima di iniziare, e nella
	\secref{sec:ciclo-ticket} descrivo il metodo con cui Halue la ottiene;
	\item a tempo e materiali, per i lavori il cui perimetro si chiarisce strada facendo. Il
	cliente paga il tempo impiegato, calcolato sul costo giornaliero di ciascun dipendente
	moltiplicato per il numero dei dipendenti coinvolti e per le giornate di lavoro;
	\item i contratti di manutenzione e assistenza per i negozi già in produzione, che coprono le
	segnalazioni e le evoluzioni dopo la consegna.
\end{itemize}

Le modifiche chieste a progetto avviato diventano \emph{change request} e seguono lo stesso
percorso di stima del progetto iniziale. Il rapporto con un cliente prosegue quindi oltre la
consegna, e con alcuni clienti è diventato una collaborazione duratura, comprensiva di più progetti
successivi alla prima collaborazione con l’azienda.

Il fondatore e il \emph{Senior Project Manager} sono le figure che seguono i clienti con più
continuità. Quando un cliente collabora da tempo con uno sviluppatore, per esempio un \emph{Senior
Software Engineer}, le richieste di nuove funzionalità e gli aggiornamenti sull’avanzamento passano
direttamente da quella persona, nelle \emph{call} e nel canale Slack dedicato al marchio.
